\chapter{Mohs Surgery}

\section*{Overview}
Mohs surgery is a tissue-sparing treatment for selected skin cancers. The visible tumor is removed in stages, and each layer is examined microscopically until no cancer cells remain at the edges.

\section*{Why It Is Done}
It is commonly used for basal cell carcinoma and squamous cell carcinoma, especially when the cancer is on the face, hands, feet, or other areas where preserving healthy tissue is important, or when the tumor has a high risk of recurrence.

\section*{How It Is Performed}
The area is numbed with local anesthetic. The surgeon removes the visible tumor and a thin surrounding layer, maps and examines the tissue, and removes additional small layers only where cancer remains. The wound is then allowed to heal or is closed with stitches or a reconstruction.

\section*{Preparation and Recovery}
Tell the clinician about blood thinners, allergies, and medical conditions. Most patients return home the same day. Wound care, bleeding precautions, and follow-up depend on the size and location of the wound.

\section*{Risks}
Bleeding, infection, pain, scarring, numbness, and recurrence are possible. The risks vary with the cancer and the repair required.

\section*{Results and Interpretation}
The procedure is complete when the examined margins contain no remaining cancer cells. Continued skin examinations are important because a history of skin cancer increases the risk of developing another lesion.

\section*{References}
\begin{enumerate}
  \item National Cancer Institute. Mohs surgery. 2025.
  \item Current Medical Diagnosis and Treatment. McGraw-Hill; 2025.
\end{enumerate}
\clearpage
